\documentclass[11pt]{article}
\usepackage{amsmath,amssymb,amsthm}
\usepackage{mathtools}
\usepackage[margin=1in]{geometry}
\usepackage{booktabs}
\usepackage{array}
\usepackage{xcolor}
\newcommand{\rsz}{\operatorname{rsize}}
\newcommand{\card}{\operatorname{card}}
\newcommand{\dlf}{\operatorname{dl}}
\newcommand{\rf}{\operatorname{rf}}
\newcommand{\acc}{\operatorname{acc}}
\newcommand{\Sstrong}{S}
\newcommand{\sfour}{\sigma_4}
\newcommand{\sseven}{\sigma_7}
\newcommand{\Aacc}{A}
\newcommand{\Bbase}{B}
\newcommand{\Dexc}{D}
\newcommand{\Uuniv}{U}
\newcommand{\sr}{\operatorname{sr}}
\newcommand{\st}{\operatorname{st}}
\newcommand{\one}{\mathbf{1}}
\newcommand{\zero}{\mathbf{0}}
\newcommand{\code}[1]{\texttt{\small #1}}
\setlength{\parindent}{0pt}\setlength{\parskip}{4pt}

\title{Route-1 Cubic Size Bound:\\How the Four Lanes Fit Together}
\author{Connecting map — COVER / BND / SEQ / FORMALIZE}
\date{2026-06-21}
\begin{document}
\maketitle

\section*{0.\quad Orientation}
There are \emph{two independent} results about the strong simplifier $\Sstrong=\code{rsimpStrong\_raw}$:
\textbf{(i) correctness} (the strong lexer returns the POSIX value) and
\textbf{(ii) the cubic size bound}. They are proved \emph{separately}; (i) may inspire (ii)'s design but is
\emph{not} assumed by it. \textbf{This note maps only track (ii)} and shows where each of the four Codex
lanes sits inside its decomposition. Every box below names the exact Isabelle lemma so each lane's own
summary can be located in the whole.

\medskip
\textbf{Primitives (self-contained).} All equations are the verbatim Isabelle definitions, in symbols.
Regexes $r,t,k$ over $\Sigma$ on the clean fragment:
\[
r\;::=\;\zero\mid\one\mid c\,(c\in\Sigma)\mid r\!\cdot\! t\mid \code{RALTS}\,rs\ (rs\text{ a list})\mid r^{*},
\]
with $\rsz\zero=\rsz\one=\rsz c=1$, $\ \rsz(\code{RALTS}\,rs)=1+\sum_i\rsz(rs_i)$, $\ \rsz(r\!\cdot\! t)=1+\rsz r+\rsz t$,
$\ \rsz(r^{*})=1+\rsz r$, and set-size $\rsz_{\mathrm s}(V)=\sum_{q\in V}\rsz q$. $\code{apder\_nf}$ = structural
normal form (no $\zero/\one$ factor in a sequence, no nested/redundant alternation); $\code{apder\_clean}=
\code{apder\_nf}$ on this fragment — the regime one may assume.

\textbf{Frontier} $\rf$ (split a top alternation, else a singleton):
\[
\rf(\zero)=\varnothing,\qquad \rf(\code{RALTS}\,rs)=\textstyle\bigcup_i\rf(rs_i),\qquad \rf(r)=\{r\}\ \text{otherwise}.
\]
\textbf{Plugs} $\sfour$ (weak), $\sseven$ (strong): $\sfour(r,k)$ appends continuation $k$ after the sequence-atom $r$,
right-associating and eliding units/zeros,
\[
\sfour(\one,k)=k,\quad \sfour(\zero,k)=\zero,\quad \sfour(r_1\!\cdot\! r_2,k)=\sfour(r_1,\sfour(r_2,k)),\quad
\sfour(h,\one)=h,\ \ \sfour(h,k)=h\!\cdot\! k
\]
for $h$ a char/$\code{RALTS}$/${}^{*}$-headed atom (and $\sfour(h,\zero)=\zero$). The strong plug adds one collapse:
\[
\sseven(a^{*},a^{*})=a^{*},\quad \sseven(a^{*},a^{*}\!\cdot\! k)=a^{*}\!\cdot\! k,\quad \sseven=\sfour\ \text{otherwise}.
\]
\textbf{Strong normaliser} $\Sstrong=\code{rsimpStrong\_raw}$ (recursive); $\Sstrong\,\zero=\zero,\ \Sstrong\,\one=\one,\ \Sstrong\,c=c$, and
\[
\Sstrong(r_1\!\cdot\! r_2)=\sseven(\Sstrong r_1,\Sstrong r_2),\quad
\Sstrong(\code{RALTS}\,rs)=\mathrm{pr}\big(\mathrm{flat}(\langle\Sstrong\,rs_i\rangle_i)\big),\quad
\Sstrong(r^{*})=\begin{cases}\one,&\Sstrong r\in\{\zero,\one\}\\ (\Sstrong r)^{*},&\text{else}\end{cases}
\]
where $\mathrm{pr}$ is the \emph{set-level cross-row prune}: in a later row $(\code{RALTS}\,lrs)\!\cdot\! k$ it deletes the
head-branches already present in an earlier row with the \emph{same} $k$, re-plugging by $\sseven$. $\mathrm{pr}$ only
\emph{shrinks} rows; its leftover ``$s^{*}\!\cdot\! s^{*}$'' survivors (uncollapsed, since $\sseven$ collapses only a
\emph{leading} equal-star) are the recurring wall.

\textbf{Opening} $\dlf=\code{row\_dlforms}$ (a regex $\to$ its set of linear rows):
\[
\dlf(\zero)=\varnothing,\ \ \dlf(\code{RALTS}\,rs)=\textstyle\bigcup_i\dlf(rs_i),\ \
\dlf\big((\code{RALTS}\,ps)\!\cdot\! k\big)=\textstyle\bigcup_p\dlf(\sseven(p,k)),\ \ \dlf(r)=\rf(r)\ \text{else.}
\]
\textbf{Strong opening of a set} $\ C(V)=\bigcup_{p\in V}\dlf(\Sstrong\,p)\ $ ($\Sstrong$ \emph{inside} the union
$\Rightarrow$ $C$ distributes over $V$).

\textbf{Term-frontier accumulator} $\acc=\code{apder\_term\_frontier\_acc}$ (continuation-parametric):
\[
\acc(\zero,k)=\acc(\one,k)=\varnothing,\quad \acc(c,k)=\rf(k),\quad \acc(\code{RALTS}\,rs,k)=\textstyle\bigcup_i\acc(rs_i,k),
\]
\[
\acc(r_1\!\cdot\! r_2,k)=\acc(r_1,\sfour(r_2,k))\cup\acc(r_2,k),\qquad \acc(r^{*},k)=\acc\big(r,\sfour(r^{*},k)\big).
\]
\textbf{Antimirov rows} $\ \code{apder\_rows}\,r=\{r\}\cup\rf(r)\cup\bigcup_{q\in\,\code{apder\_terms}\,r}\rf(q)\ $ (the
deduped one-step row index set), and the \textbf{target universe} $\Uuniv(r)=C(\code{apder\_rows}\,r)$.

\medskip
\textbf{Carrier, base, singletons, excess (the single letters used below).}
\[
\Aacc(r,k)=C\big(\rf(\sfour(r,k))\cup\acc(r,k)\big)\ \ (=\code{strong\_apder\_acc}),\qquad \Bbase(k)=\Aacc(\one,k),
\]
\[
\sr(q,k)=C\big(\rf(\sfour(\code{RALTS}[q],k))\big),\quad \st(q,k)=C(\acc(q,k)),\quad \Aacc(\code{RALTS}[q],k)=\sr(q,k)\cup\st(q,k),
\]
\[
\Dexc(r,k)=\card\big(\Aacc(r,k)\setminus\Bbase(k)\big),\qquad \Dexc_1(q,k)=\Dexc(\code{RALTS}[q],k).
\]
The final-theorem wrappers $\code{afactored1}$, $\code{rpder\_strong\_rows\_raw}$, $\code{row\_dlformss}$ are
\emph{opaque}: they occur only inside the green gate black-box (\S3) and are used nowhere else here.

\section*{1.\quad The final theorem (the top)}
\[
\rsz\!\big(\code{row\_dlformss}\ (\code{rpder\_strong\_rows\_raw}\ c\ (\code{afactored1}\ r\ s))\big)\;\le\;2\,(\rsz r+3)^3 .
\]
Plain words: open every one-step strong-pruned derivative row, dedupe, and the total size is cubic in $\rsz r$.

\section*{2.\quad The gate: reduce ``cubic'' to one \emph{linear} count}
\code{DirectUniverseCubic.thy} is \emph{green} except for one obligation. With the universe
$\Uuniv(r)=\code{apder\_strong\_dlfrontier}\ r=C(\code{apder\_rows}\ r)$, the proved gate lemmas
(\code{actual\_gate\_from\_direct\_universe\_rowlevel}, \code{universe\_le\_cubic\_rowlevel},
\code{per\_row\_size\_le\_quadratic}) give the final theorem \emph{provided}
\[
\boxed{\ \code{card\_apder\_strong\_dlfrontier\_le}:\quad \card\,\Uuniv(r)\ \le\ \rsz r+1\ }\tag{G}
\]
(each member has \emph{quadratic} size; cubic $=$ quadratic size $\times$ \emph{linear} count). Everything below
is the proof of the single linear-count obligation (G). \textbf{Any} linear bound $C\!\cdot\!\rsz r+D$ also closes it.

\section*{3.\quad Bridge: from $\Uuniv$ to the excess counter $\Dexc$}
Three green facts collapse (G) to a statement about $\Dexc$ at continuation $k=\one$:
\[
\Uuniv(r)\subseteq \Aacc(r,\one)\ \ (\code{\dots\_subset\_strong\_apder\_acc\_RONE}),\quad
\Aacc(\one,\one)=\{\one\}\ (\code{strong\_apder\_acc\_RONE\_RONE}),
\]
\[
\Rightarrow\quad \card\,\Uuniv(r)\ \le\ \card\,\Aacc(r,\one)\ \le\ 1+\Dexc(r,\one)
\quad(\code{card\_le\_Suc\_card\_Diff\_singleton}).
\]
So (G) follows from the \textbf{global excess bound}
\[
\boxed{\ \code{card\_strong\_apder\_acc\_diff\_base\_le\_rsize}:\quad \Dexc(r,k)\ \le\ \rsz r\quad\text{(all }k\text{)}\ }\tag{$\Delta$}
\]

\section*{4.\quad The decomposition (how $(\Delta)$ splits down to the four lanes)}
$(\Delta)$ is proved by induction on $r$ (arbitrary $k$):
\begin{itemize}
\item $\code{RZERO}/\code{RONE}/\code{RCHAR}$: base cases ($\Dexc\le \rsz$ directly).
\item $\code{RSEQ}\,r_1 r_2$: the \textbf{RSEQ recurrence} (below).
\item $\code{RALTS}\,rs$: the \textbf{RALTS budget} $\ \Dexc(\code{RALTS}\,rs,k)\le \sum_{q\in rs}\rsz q\ $
      (\code{\dots\_size\_budget\_spine}), which needs \textbf{L1} $+$ \textbf{L2}.
\item $\code{RSTAR}\,r$: two star helpers (root-row $\le 1$; the $\sfour(\code{RSTAR})$ frontier identity).
\end{itemize}

\textbf{L2 (singleton size).} $\ \boxed{\Dexc_1(q,k)\le \rsz q}\ $ (\code{card\_strong\_apder\_acc\_singleton\_RALTS\_diff\_base\_le\_rsize}),
by induction on $q$. Its $\code{RSEQ}$ case is exactly the
\[
\textbf{RSEQ recurrence:}\qquad \Dexc_1(\code{RSEQ}\,r_1 r_2,\,k)\ \le\ \rsz r_1+\Dexc_1(r_2,k)
\quad(\code{\dots\_RSEQ\_rec\_spine}),
\]
and \emph{that} recurrence is what consumes the two boundary lemmas \textbf{S1} and \textbf{seq\_head}. Its
$\code{RALTS}$ case re-uses \textbf{L1}.

\medskip
\textbf{The three leaf obligations the lanes own:}
\[
\begin{aligned}
\textbf{L1 (cover):}\quad & \Aacc(\code{RALTS}\,rs,k)\ \subseteq\ \textstyle\bigcup_{q\in rs}\Aacc(\code{RALTS}[q],k).\\[2pt]
\textbf{S1 (boundary absorb):}\quad & \card\!\big((\Aacc(\one,\sfour(t,k))\!\setminus\!\Bbase(k))\cup(\st(t,k)\!\setminus\!\Bbase(k))\big)\le \Dexc_1(t,k).\\[2pt]
\textbf{seq\_head:}\quad & \card\!\big((\sr(\code{RSEQ}\,h\,t,k)\cup\st(h,\sfour(t,k)))\setminus(\Bbase(k)\cup\Bbase(\sfour(t,k)))\big)\le \rsz h.
\end{aligned}
\]

\medskip
\textbf{Top-to-leaf chain, one line:}
\[
\text{cubic}\ \xLeftarrow{\;\text{gate}\;}\ (\mathrm{G})\ \xLeftarrow{\;\text{bridge}\;}\ (\Delta)\ \xLeftarrow{\;\text{ind. }r\;}\ \{\text{RSEQ rec},\ \text{RALTS budget},\ \text{RSTAR}\}\ \xLeftarrow{}\ \{\textbf{L1},\ \textbf{L2}\}\ \xLeftarrow{\;\text{ind. }q\;}\ \{\textbf{S1},\ \textbf{seq\_head},\ \textbf{L1}\}.
\]

\section*{5.\quad The four lanes, placed in the map}
\renewcommand{\arraystretch}{1.35}
\begin{center}\small
\begin{tabular}{@{}>{\bfseries}l p{0.30\linewidth} p{0.40\linewidth}@{}}
\toprule
Lane & Proves (Isabelle) & Role / method \\
\midrule
COVER & \code{strong\_apder\_acc\_RALTS\_singleton\_cover} (\textbf{L1}) &
Feeds the \emph{RALTS budget}. Method: fix-(a), two carriers; the cross-prune-collapsed $s^{*}$ escapes are
recovered by the \emph{acc} carrier, not a branch root. \\
BND & \code{boundary\_term\_absorb} via \textbf{S1} (\emph{OPEN}) &
Feeds the \emph{RSEQ recurrence}. The card target $\card X\le\card Y$ is OPEN: the asymmetric
$\Sstrong$-image repair was REFUTED (ad485d1) even for RSEQ-root; a missing-image repair
$X\setminus Y\subseteq\Sstrong\text{-image}(Y\setminus X)$ is proposed but unvalidated. \\
SEQ & \code{seq\_head\_core\_le\_rsize} (\textbf{seq\_head}) &
Feeds the \emph{RSEQ recurrence}. Method: charge $\rsz h$; \code{single\_root\_RSEQ\_RCHAR\_card\_le\_one}
helper; induction on $h$ (uses \textbf{L1} on the RALTS-head case). \\
FORMALIZE & the \emph{spine}: bridge, $(\Delta)$ induction, RSEQ-rec assembly, RALTS budget, RSTAR helpers, $\Rightarrow$ (G) &
Assembles \S3--\S4 on \code{assumes} for L1/S1/seq\_head, swapping each in as its lane lands it. \\
\bottomrule
\end{tabular}
\end{center}

\section*{6.\quad Dependencies / merge order}
\code{COVER}(L1) and \code{BND}(S1) are independent leaves $\Rightarrow$ land in parallel first.
\code{SEQ} depends on \emph{both} (its RSEQ step uses S1; its RALTS-head case uses L1).
\code{FORMALIZE} integrates all three; once L1, S1, seq\_head are green it discharges every \code{assumes},
proves $(\Delta)\Rightarrow$(G)$\Rightarrow$ the cubic theorem. Status (track ii): the spine and many helpers are
\emph{green}; L1 / S1 / seq\_head are \emph{validated with proof skeletons}, in formalization.
\end{document}
